\documentclass{article}
\usepackage{amsmath}
\usepackage{amssymb}
\usepackage{hyperref}
\usepackage{url}

\title{Compliant Grasping and Causal Frame Selection:\\A Conditional Boundary Between Aero Hand Open and DynaMAC}
\author{}
\date{}

\begin{document}
\maketitle

\begin{abstract}
This note relates Aero Hand Open, a compliant tendon-driven robot hand, to DynaMAC, a method for coordinating manipulation through moving reference frames. We asked whether DynaMAC's test for a rigid link can identify causal attachment when a compliant grasp allows relative motion. The papers together make a possible mismatch clear, and the peers designed a controlled test of it, but no experiment was run. The thread therefore paused because the workspace lacked an integrated simulator and an agreed definition of attachment.
\end{abstract}

\section{Introduction}

Aero Hand Open is an anthropomorphic robot hand with 16 joints and seven motors \cite{wang2026aero}. Cables let one motor move several joints. Springs and silicone contacts make the hand compliant, so the mechanics and the object help determine which joints move during closure. The hand senses its seven motor encoders. Its released learning example concerns fixed-base rotation of a cube, and its policy sends seven tendon targets. The paper also provides a simulation model, an identified map between motor and simulated commands, and a deployment stack.

DynaMAC addresses dynamic and two-arm manipulation \cite{hartz2026dynamac}. It represents actions relative to several reference frames, such as an object or the other arm's end effector. Each stream predicts an absolute end-effector pose, and streams with lower predictive uncertainty receive greater weight. DynaMAC adds a binary test for a kinematic link. When a frame appears rigidly attached, the method masks the stream whose frame has become an effect of the robot's action. DynaMAC disregards gripper actions, and DynaBench tests arm end-effector motion rather than dexterous hand control.

The connexion pursued was narrower than mounting Aero hands on a DynaBench system. A grasp can be causal in the sense that the hand controls the object, yet non-rigid in the sense that the object can slip or move within compliant contacts. The peers asked whether DynaMAC's rigid-motion test then selects the right streams. They also asked whether a test based on the object's response to small wrist and tendon interventions would do better. This is a question about identifying attachment, not a claim that Aero makes DynaMAC fail.

\section{What Was Done}

The work began with a broad proposal to replace standard grippers with Aero hands. The peers rejected a direct comparison as ambiguous. DynaMAC controls a six-dimensional arm pose and omits gripper actions, whereas the Aero example controls seven tendon targets in a fixed-base task. A failed combined system could therefore reflect hand control, mounting, sensing, collision modelling, or task design rather than causal frame selection.

The peers then isolated the common issue. The arm's wrist pose remains known from forward kinematics. What the wrist pose does not fully describe is the position and strength of compliant finger contacts. Aero is a useful counter-regime because a cable can drive several joints, the mechanics choose the order of joint motion, and contact pads deform. Thus a known wrist pose and tendon command need not determine a rigid object pose.

Next, the peers examined DynaMAC's link rule. Let $M$ be the geometric mean of the standard deviations of relative pose, including translation and rotation. DynaMAC declares a link when
\begin{equation}
  M < \tau_M,
\end{equation}
where $\tau_M=0.001$. The paper interprets this threshold as about 1~mm or 0.001~rad. It is an operational test motivated by rigid grasping. DynaMAC also notes brief false positives before a grasp.

An early proposal replaced this binary decision with a continuous covariance weight. The peers withdrew it. DynaMAC's underlying product of experts already weights streams by predictive precision, so another covariance weight would repeat an existing mechanism. More importantly, variance alone cannot distinguish an external object from an attached object that slips or makes intermittent contact.

The revised proposal separated attachment from predictive precision. Let $z_t$ denote whether the object is attached at time $t$. The proposed estimator would apply matched, small wrist and tendon interventions and ask whether the object or contact frame responds. It would use the estimated probability that $z_t$ is attached to gate a stream, while leaving ordinary covariance weighting unchanged. This was a design only. No estimator, label, or experimental result was produced.

Finally, the peers specified a simulation-first comparison. Grasp control would be frozen or identical for every method. A model would vary from rigid coupling through compliant slip to intermittent contact. The comparison would include ordinary precision weighting, binary DynaMAC, a gate given an independently defined attachment label, and a learned response-based gate. The two main outcomes would be agreement with the attachment label and success when policies learned from static demonstrations face dynamic perturbations.

\section{Results}

Three careful conclusions hold on the supplied evidence.

First, Aero provides a plausible setting in which causal control and rigid co-motion separate. Its underactuation, mechanically staged closure, compliant contacts, and limited actor observations support that possibility. They do not establish that the separation occurs in DynaBench or that it changes task success.

Second, DynaMAC uses near-constant relative pose as a practical sign of a kinematic link. Therefore its binary proxy may miss an attachment that transmits control while allowing relative motion. It may also give a false positive during intermittent or pre-grasp contact. These are conditional failure modes, not a proved defect. Compliance could instead keep a stream's uncertainty finite and prevent the excessive precision that masking was introduced to correct.

Third, response-based attachment inference is conceptually different from DynaMAC's variance test and from ordinary precision weighting. It asks whether a controlled change in the hand produces a change in the object. The peers' cross-checks support this distinction, but only as an experimental hypothesis. The supplied \texttt{ada/} and \texttt{emmy/} directories contain no additional files, so the ledger is the only available record of those checks.

The proposed experiment has clear interpretations. If the binary rule and the independent attachment label agree, the concern is refuted in the tested regime. If they disagree but task success is unchanged, the distinction has little demonstrated consequence. If they disagree and the independently labelled or learned gate improves transfer, the proposed boundary becomes consequential.

\section{What Did Not Hold}

The broad hardware-swap claim did not hold up. A direct replacement would mix arm coordination with a new seven-channel hand controller. Holding grasp control fixed, or giving every method the same tendon controller, settles that objection for the proposed causal-frame test. It does not establish dexterous two-arm policy learning or DynaMAC's sample efficiency with Aero hands.

The continuous covariance proposal also did not hold up because ordinary product-of-experts weighting already uses predictive covariance. The response-based variable $z_t$ removed that duplication in principle. It did not settle how attachment should be defined during intermittent contact involving several fingers.

The peers also corrected an overstatement about observability. The wrist is observable. The missing information concerns contact geometry and force transmission. On hardware, a response-based test would need external object or contact tracking because Aero exposes only motor encoders.

No evidence showed that DynaMAC's threshold disagrees with an independent attachment label or that any disagreement harms transfer. Neither paper performs the proposed combined test. No literature search was conducted, so the thread supports no claim that the proposal is new beyond these two papers.

\section{Why the Thread Stopped}

After the first review, the verifier requested the rigid-to-compliant ablation. The peers then checked the workspace and found no Aero or DynaMAC codebase, demonstrations, trained policies, task assets, or integrated Aero--DynaBench environment. The Aero project points to a public code repository \cite{aerocode}, and DynaMAC points to a project release \cite{dynamacrelease}, but those links do not supply a ready combined experiment in this workspace. The verifier therefore changed the status from \textsc{iterate} to \textsc{pause}: the idea was testable, but the decisive evidence could not be produced from the supplied material.

The smallest step that would restart the thread is to obtain the released code and build a single integrated simulation case with fixed grasp control. Before running it, the researchers must define an attachment label for rigid, slipping, and intermittent contact. They can then compare binary DynaMAC with the independently labelled gate on label agreement and transfer success. Testing the learned response-based gate is needed only if that first comparison shows a consequential disagreement.

\bibliographystyle{plain}
\bibliography{references}

\end{document}
